\section{Limitations and Future Work}
\label{sec:limitations}

While our framework provides a verifiable approach to generating discrete-event world models, several limitations define its current scope and present directions for future research:

\textbf{Scope of Environment Dynamics.} As motivated, \methodname{} is purposely designed for macro-level, discrete-event environments (e.g., supply chains, service operations, and logical workflows). Consequently, it does not natively model continuous-time differential equations. While this strict discrete scope is sufficient for process-driven domains, some advanced real-world applications—such as smart grids or automated cyber-physical manufacturing—exhibit \emph{hybrid} dynamics, where discrete strategic events intersect with continuous physical evolution. Extending the framework to synthesize Hybrid DEVS models to support these intersecting continuous domains remains a technical direction for future work.

\textbf{Instruction Adherence and Model Fine-Tuning.} The reliability of \methodname{}, particularly on smaller models, is bounded by the underlying LLM's instruction-following capabilities. Generation failures frequently stem from minor deviations in output formatting or port naming rather than fundamental logic errors. Because our decoupled workflow decomposes the generation process into standardized, predictable sub-tasks (\modelplan{}), the prompt distribution is highly stable. Future work can leverage this property by applying Supervised Fine-Tuning (SFT) or Reinforcement Learning from Human Feedback (RLHF) to smaller coding models, which is expected to significantly mitigate formatting failures.

\textbf{Benchmark Scale and Specification Ambiguity.} We introduce a structured benchmark to evaluate simulator synthesis; however, the current dataset relies on a bounded set of scenarios with precise, unambiguous system specifications. This controlled setup is necessary to isolate pure synthesis capability from requirement analysis. In practical deployments, target systems can scale to thousands of interacting components, and user prompts are often incomplete. Future iterations will expand the benchmark to include larger-scale, highly underspecified tasks. Correspondingly, the framework can be enhanced with an interactive elicitation phase to clarify ambiguous requirements, alongside a targeted self-correction loop in Stage 1 to resolve structural topology errors before code synthesis.